\subsection{Iterating the actual leading-minor maps}
\label{archimedean-iterated-minors}

The local family just constructed also supplies the original
higher weighted norm estimates. The pointwise orbital bound remains
a separate assertion: all equalities between orbital integrals in
this subsection are equalities in the original radial $L^2$ space.

Write $q=r+1$, $m=2q$, and $n=2r+k+1$, with $k\ge1$.
Retain the original local datum, section, character scale
$\kappa\in\{1,a\}$ and radial measure. Put
\[
 Q=q_0(C_q,m,\epsilon,u,\lambda_0),\qquad
 \delta=\lambda-\lambda_0\in\mathbb Z^n .
\]
For $\alpha=(\alpha_1,\ldots,\alpha_q)\in\mathbb N^q$ put
$|\alpha|=\sum_{\ell=1}^q\alpha_\ell$ and
$D^\alpha=\prod_{\ell=1}^qD_\ell^{\alpha_\ell}$, with the
original leading minors $D_\ell$.

\begin{theorem}\label{original-iterated-minor-norm}
Assume the original $\mathrm{LocOk}_0$ and near-parameter range
hypotheses, including
\[
 \|\delta\|_1+q(p+q+1)\le Q,\qquad
 \text{every gap of }\lambda_0\text{ is at least }C_gp\,m .
\]
If $C_q\ge200$, $C_g\ge10^6$ and $C_b\ge C_h+5$, then, for
every original compact-finite coefficient $F$ and every
$|\alpha|\le p$,
\[
 \mathcal N_\kappa(D^\alpha F)
 \le b^{\,2C_bm^2|\alpha|}\mathcal N_\kappa(F).
\]
This supplies both the original near-parameter iteration statement
and, on setting $\lambda_0=\lambda$, its derivative-iteration
statement \cite{iterated-minors-source}.
\end{theorem}

\begin{proof}
Every norm in the proof uses the same original radial measure.
The integral lattice and flooring argument of
Proposition~\ref{jacquet-original-budget-transport} gives the
same family at all integer parameters needed below.
For $\eta\in\mathbb N^n$ with $\|\eta\|_1\le qp$, the parameter
$\lambda+\eta$ lies in this family, and
\[
 \|\delta+\eta\|_1+q(q+1)
 \le\|\delta\|_1+qp+q(q+1)\le Q.
\]
Thus the full leading-minor reserve remains available at every
intermediate parameter. The original all-gap hypothesis is retained.
No numerical constant changes with $\eta$ or the chosen word.

The mixed Jacquet identity says that the images of the coefficients
of one representation form a subspace of dimension at most one in
this radial $L^2$ space. The nonzero such subspaces at distinct
parameters are orthogonal. Choose a unit vector $\Phi_\eta$ in each
nonzero subspace; choose the zero vector when that subspace is zero.
These choices can be represented by original compact-finite
coefficients: a nonzero normalized Bessel map is nonzero on
compact-finite vectors, and a rank-one operator between its two
nonzero images supplies a coefficient with nonzero mixed period.
The common norm identity then normalizes it.
No nonvanishing is asserted for a subspace whose period is zero.

For $1\le\ell\le q$, the actual exterior-Pieri decomposition gives
\[
 D_\ell\Phi_\eta
    =\sum_{|I|=\ell}a_{\eta,I}^{(\ell)}
                          \Phi_{\eta+\mathbf1_I}
\]
in this $L^2$ space. A term with zero target subspace is zero.
If the source subspace is zero, set all coefficients to zero.
For a nonzero source and target, the exact Gram identity gives
\[
 |a_{\eta,I}^{(\ell)}|^2
   =|\kappa_{\eta,I}\kappa'_{\eta,I}|^2
                              \frac{d_{\lambda+\eta}}
                                   {d_{\lambda+\eta+\mathbf1_I}} .
\]
The period-normalized moment bounds therefore imply
\[
 |a_{\eta,I}^{(\ell)}|^2
 \le \kappa^{4S_\ell}(TT')^\ell b^{8m^2}
             \frac{d_{\lambda+\eta+\mathbf1_I}}
                  {d_{\lambda+\eta}},\qquad
 S_\ell=r\ell-\ell(\ell-1)/2.
\]
Apply Cauchy--Schwarz and the exact degree-ratio sum:
\[
 \sum_{|I|=\ell}|a_{\eta,I}^{(\ell)}|
 \le \binom n\ell\,\kappa^{2S_\ell}(TT')^{\ell/2}b^{4m^2}.
\]
Proposition~\ref{jacquet-original-scale-transport} gives
$\kappa^{2S_\ell}(TT')^{\ell/2}\le b^{C_hm^2}$.
Also $\binom n\ell\le n^\ell\le b^\ell$ and
$\ell\le q\le m^2$. Hence every such column obeys
\[
             \sum_I|a_{\eta,I}^{(\ell)}|
                       \le B:=b^{C_bm^2}.
\]
This sum bound is needed because different paths can reach the
same final parameter.

Choose any ordered list of the $|\alpha|$ leading minors in
$D^\alpha$. Starting with a unit vector at $\eta=0$, expand
one factor at a time. If its current coefficients are $c_\eta$,
the sum of the absolute values after the next factor is at most
$B\sum_\eta|c_\eta|$. Combining terms at coincident endpoints
uses the triangle inequality. Induction therefore gives
\[
 D^\alpha\Phi_0=\sum_\eta c_\eta\Phi_\eta,\qquad
 \sum_\eta|c_\eta|\le B^{|\alpha|},\qquad
 \|\eta\|_1\le q|\alpha|\quad(c_\eta\ne0).
\]
All sums have finite support. Orthogonality now yields
$\|D^\alpha\Phi_0\|_2^2
\le(\sum_\eta|c_\eta|)^2\le B^{2|\alpha|}$.

The original norm of $F$ is finite by the mixed norm identity.
If it is positive, choose
$\Phi_0=J_\kappa F/\sqrt{\mathcal N_\kappa(F)}$ and rescale
the preceding inequality. If it is zero, $J_\kappa F=0$
almost everywhere, and multiplication by the finite polynomial
$D^\alpha$ gives the conclusion directly. The empty word also
has the asserted estimate.
\end{proof}

For any fixed admissible datum, finite shift range and chosen
coefficient representatives, the finitely many expansions can be
realized off one common null set for the original radial measure.
Its choice may depend on that datum and those representatives.
The argument neither chooses one null set for every arithmetic
parameter nor changes a statement about a specified rational point
into an almost-everywhere statement.

In particular it proves the norm consumers
$\mathrm{BesselIterationNearParameter}$ and
$\mathrm{BesselDerivativeIteration}$ exactly.
It does not yet prove the stronger source coefficient-expansion
interface, which requires equality at every big-cell point.
Nor does it prove the weighted orbital bound on every relevant
stratum. The final rational-return sum can be supported on a
Haar-null set, so that pointwise frontier cannot be removed
by Tonelli alone.
